\section{Support evidence}
\label{sec:support}

Nothing here is a further finding. This section holds what the three findings
rest on: the semantics of the two governing predicates read at source
(Section~\ref{sec:impl}), an execution of the published relations on a real
system (Section~\ref{sec:native-run}), and the robustness of our numbers to a
second, independent implementation of the same predicates
(Section~\ref{sec:secondimpl}); the seeded positive control
(\S\ref{sec:seeded}) and the cross-vendor replication (\S\ref{sec:gogs}) are the
remaining confirmations and stay with the study that produced them. Each item is
a confirmation of a mechanism already stated, not a new claim about what the
technique can detect.

\subsection{Confirmation at source}
\label{sec:impl}

Definitions can be read generously, so we checked the predicates at source in
the public library~\cite{Chaleshtari2022Metamorphic}. It ships \ResArtMrTotal{}
relations, \ResArtAuthzMrTotal{} authorization relations,
\ResArtAuthzGatedEither{} gating on the two predicates analysed. Four
implementation-level properties not made explicit in the main presentation:
(i)~\emph{the freeze is hard-coded}---%
\texttt{cannotReachThroughGUI} is set-membership over the crawl inputs, issues
no request, and is cached in a field written once and never invalidated
(\texttt{userCanRetrieveContent} likewise reads a crawl-produced store loaded
once); (ii)~\emph{the clause that could have caught it is inert}---the
``absent across all users'' set is built by stopping each recorded sequence at
the first login action, so it is always empty; (iii)~\emph{missing data defaults
to an alarm}---with an empty or null crawl argument the predicate returns
``cannot reach'', so an empty crawl \emph{satisfies} the precondition: the
mechanism behind Finding~1's deterministic direction; (iv)~\emph{the role
relation is configuration, not state}---the supervisor guard reads a static
configuration file and reports true for identical accounts, so it cannot reflect
a later grant.

\subsection{Executing the published relations}
\label{sec:native-run}

The library ships relations but no execution environment; we supplied the system
configuration and crawl snapshot and ran the relations---source unmodified,
build configuration (\texttt{pom.xml}) and an external harness added
(\NtUpstreamModified{} changed upstream file, \NtUpstreamJavaModified{} changed
Java sources, two added classes totalling \NtHarnessLines{} lines)---against
Gitea \NtGiteaVersion{} (Table~\ref{tab:native-facts}). \NtMrExecuted{} of \NtMrTotal{} authorization relations
executed; \NtMrFired{} returned failures, \NtAlarms{} alarms in all---the
engine's own \texttt{getFailures()} value (Table~\ref{tab:native}). What the run
establishes is scoped to the \emph{reachability-gated subset that fires here},
not a full replay of the catalogue: one relation aborts on an engine-side defect,
most issue no request because their own guards are unsatisfied, and the content
predicate is never invoked. Almost none
of the \NtElapsedS{} seconds is analysis: each leaf evaluation invokes a
browser-cleanup hook costing minutes, a cost of the missing environment that
bounds any reproduction.

\begin{table}[t]
\centering
\caption{Execution facts for the native run of the published relations
(Section~\ref{sec:native-run}). All values are emitted by
\texttt{ingest\_native\_engine.py} from the run record; upstream Java sources are
unchanged apart from the build configuration.}
\label{tab:native-facts}
\small
\begin{tabular}{@{}p{0.66\columnwidth}r@{}}
\toprule
item & value \\
\midrule
authorization relations: catalogue / executed / fired & \NtMrTotal{} / \NtMrExecuted{} / \NtMrFired{} \\
engine alarms (the engine's own return value) & \NtAlarms{} \\
HTTP requests attributed to observing relations & \NtHttpCalls{} \\
changed upstream Java sources & \NtUpstreamJavaModified{} \\
build/config changes (one \texttt{pom.xml} line) & \NtUpstreamBuildModified{} \\
added harness classes / added lines & two / \NtHarnessLines{} \\
\bottomrule
\end{tabular}
\end{table}

The load-bearing cell: the subject \NtKeyGrantee{} holds a legitimate
entitlement to \NtKeyRepo{}, granted out of band after the crawl. There the
published \texttt{cannotReachThroughGUI} answers ``cannot reach''---the answer
it is \emph{defined} to give, since the subject's crawl record has no entry for
the repository; across the run it answered so \NtCrtgTrue{} of \NtCrtgCalls{}
invocations (\NtCrtgTruePct{}\%), \NtKeyCrtgTrue{} on this cell, and
\NtKeyMrFired{} relations (\NtKeyMrNames{}) report a failure there. The access
is authorized: all \NtByteCells{} granted cells were re-probed, and the
grantee's response hashes identically to the owner's while a third account is
refused. The \NtAlarms{} alarms fall on \NtAlarmCells{} cells, all post-crawl
grants\ifnum\NtAlarmsOffGrant=0, none on a non-entitled object.\else, plus
\NtAlarmsOffGrant{} on a non-entitled object.\fi{} Limits: the harness's HTTP
observation channel is exact for the reachability predicate (set-membership, no
request) but unexercised for the content predicate (\NtUcrcCalls{} invocations;
supervisor guard \NtSupCalls{}, \NtSupTrue{} true); not all gated relations fire,
for reasons inside their own guards, and \NtMrErrored{} abort on an engine-side
defect---so the reproduced false proof is that of the reachability-gated subset
that fires, not a full-replay count of the published relations. And Gitea, on
the tested surface, enforces the expected decisions, so the run can only exhibit
false proofs.

\begin{table*}[t]
\centering
\caption{The published authorization relations executed against Gitea
\NtGiteaVersion{}. Each row is the applicability chain from the catalogue to the
outcome: \emph{status} is the engine's per-relation result (\NtMrExecuted{} of
\NtMrTotal{} evaluated, \NtMrErrored{} aborted on an engine-side defect),
\emph{alarms} its own \texttt{getFailures()} return (\NtMrFired{} relations,
\NtAlarms{} in all), and \emph{requests} (\NtHttpCalls{} total) attributed to the
relation observing them first.}
\label{tab:native}
\small
% 由 analysis/ingest_native_engine.py 自动生成 —— 请勿手改
% 逐条 MR 结果：告警数由引擎自己 getFailures() 返回，非本工作判读
\begin{tabular}{lrr}
\toprule
MR & alarms & requests \\
\midrule
OTG\_AUTHZ\_001 & 0 & 0 \\
OTG\_AUTHZ\_001b & 0 & 0 \\
OTG\_AUTHZ\_001b2 & 0 & 0 \\
OTG\_AUTHZ\_002 & 3 & 270 \\
OTG\_AUTHZ\_002a & 0 & 0 \\
OTG\_AUTHZ\_002b & 0 & 0 \\
OTG\_AUTHZ\_002c & 3 & 0 \\
OTG\_AUTHZ\_002d & 0 & 0 \\
OTG\_AUTHZ\_002e & \multicolumn{2}{c}{error} \\
OTG\_AUTHZ\_003 & 0 & 0 \\
OTG\_AUTHZ\_004 & 0 & 0 \\
\bottomrule
\end{tabular}

\end{table*}

\subsection{Implementation robustness}
\label{sec:secondimpl}

An independently written implementation of the same predicates was compared with
the main one on a subset fixed before the comparison, not sampled after it: the
in-band scenarios on which a second reading is defined, holding the
exception-policy variant at the representative level and the reachability reading
at the one each relation uses (MR-002 blind, MR-004 aware). This yields
\ResAltPoints{} (configuration, scenario) points; MR-002 is defined on
\ResAltMrTwoPoints{} of the resulting relation comparisons and MR-004 on the
remaining \ResAltMrFourPoints{}, giving \ResAltCompared{} in all. The two agreed
on \ResAltAgree{} of them (\ResAltAgreeRate{}); the point set is a fixed
enumeration, so no interval is attached. Both were written from the same
reading of the semantics, so this bounds implementation accident only; it does
not bound a shared blind spot, and we do not present it as corroboration.
